% Tabla 1: Trazabilidad entre preguntas, objetivos, hipótesis, resultados y conclusiones.
% Etiqueta original: tab:trazabilidad.
% Fuente: capitulos/1-introduccion.tex, lineas 73-89.
% Fragmento extraido automaticamente; requiere el preambulo de main.tex.

\begin{table}[htbp]
\centering\footnotesize
\caption{Trazabilidad entre preguntas, objetivos, hipótesis, resultados y conclusiones.}
\label{tab:trazabilidad}
\begin{tabular}{@{}cp{4.6cm}cccc@{}}
\toprule
Obj. & Qué produce & Pregunta & Hipótesis & Resultado & Conclusión \\
\midrule
1 & Conjunto de datos documentado & --- & --- & §\ref{sec:conjunto} & §\ref{sec:conclusiones}, c.~1 \\
2 & Protocolo y repetibilidad del marcado & S2 & --- & §\ref{sec:procedimiento} & §\ref{sec:conclusiones}, c.~2 \\
3 & Rutas A, \Aprime, B y C implementadas y validadas & S1 & H1 & §\ref{sec:calibraciones}, §\ref{sec:agregados} & §\ref{sec:conclusiones}, c.~3 \\
4 & Comparación pareada por separación & S1 & H1, H2 & §\ref{sec:profundidad}, §\ref{sec:contraste} & §\ref{sec:conclusiones}, c.~4 \\
5 & Intervalo $\IL$, ancho y cobertura & Principal & H3 & §\ref{sec:intervalo} & §\ref{sec:conclusiones}, c.~5 \\
6 & Sensibilidad a marcado y geometría & S2, S3 & --- & §\ref{sec:sensibilidad} & §\ref{sec:conclusiones}, c.~6 \\
\bottomrule
\end{tabular}
\end{table}
